\section{Data and formal verification}
\label{app:finite-certificates}

This appendix specifies the $48$ exact positive-increment steps used in
Lemma~\ref{lem:finite-increments} and the exact tests that establish
their validity. The accompanying files
\path{verification/data/step_02.json.gz} through
\path{verification/data/step_49.json.gz} contain rational
coefficients and their indices. The manifest lists the steps in
order and gives hashes of both compressed and uncompressed data.

\subsection{The two types of increment}\label{app:increment-data}

For $2\le n\le16$, the certificate represents
\begin{equation}\label{eq:appendix-normalized-step}
 R_n=(n-1)F_{n+1}-(n+1)F_n
\end{equation}
as an SoS. In \eqref{eq:finite-chain} this gives
$\alpha_n=(n+1)/(n-1)$, the coordinate projection for $\psi_n$,
and the Gram of $R_n$ divided by $n-1$ for $\mathbf H_n$.
For $17\le n\le49$, it represents
\begin{equation}\label{eq:appendix-geometric-step}
 I_n=F_{n+1}-F_n\circ\psi_n,\qquad
 \psi_n(x,y)=
 \bigl((q_n^{|a|}x_a)_{|a|<n},(q_n^{|a|}y_a)_{|a|<n}\bigr),
\end{equation}
where $q_n=1-\theta_n/n$ is rational and $0<q_n\le1$.
The rational number $\theta_n$ is stored in the corresponding data
file. Here $\alpha_n=1$. In particular, $q_{49}=97/98$ at the final step.

The stored Gram blocks use noncentral wedges, with both labels
nonzero. Central wedges are added as separate squares. For each
$1\le a\le n$, the coefficient of each of
$z_{-a,0}^2$ and $z_{0,a}^2$ is
\[
 \begin{cases}
 2(n+1)(a-1),&\text{in }R_n,\\
 2\bigl((n+1)(n+1-a)-q_n^{2a}n(n-a)\bigr),
     &\text{in }I_n.
 \end{cases}
\]
These coefficients are nonnegative. Adjoining them produces a
Gram on all wedges, as required in \eqref{eq:finite-chain}.

\subsection{Exact polynomial and positivity checks}
\label{app:exact-psd-checks}

For each step, the checker expands $F_n$ and $F_{n+1}$ directly
from the Toeplitz matrix definition. It also expands the proposed
Gram expression and verifies equality of every monomial coefficient.
The Pl\"ucker corrections are separately checked to represent zero.
All coefficients in these comparisons are integers or rationals.

The full noncentral Gram has verified kernel vectors
\[
 h_g=\sum_{\substack{-n\le a<0<b\le n\\b-a=g}}e_{ab},
 \qquad n+1\le g\le2n,
\]
where $e_{ab}$ is the coordinate vector for $z_{ab}$.
Their supports are disjoint. For each $h_g$, delete the first
coordinate in its support, in lexicographic order. The remaining
coordinate vectors together with the $h_g$ form a basis. Once
the exact relations $\mathbf H_nh_g=0$ have been checked, this change of
basis gives a congruence
\[
 \mathbf H_n\ \sim\ \begin{pmatrix}\widehat H_n&0\\0&0\end{pmatrix},
\]
where $\widehat H_n$ is the retained principal submatrix.
Consequently $\mathbf H_n\succeq0$ follows from $\widehat H_n\succeq0$.
No theorem about the kernel of all possible Grams is needed for
this check of the supplied matrices.

Each connected block of $\widehat H_n$ is checked independently.
One available test is Sylvester's criterion using exact
fraction-free elimination. The default checker uses an integer
residual certificate. Write a rational block as $A/d$, with
$A$ integral and $d>0$. For an integer matrix $L$ and integer
$s\ge0$, it computes
\begin{equation}\label{eq:exact-residual-certificate}
 E=2^{2s}A-dLL^\top,\qquad
 \delta=\min_i\left(E_{ii}-\sum_{j\ne i}|E_{ij}|\right).
\end{equation}
If $\delta>0$, symmetry and diagonal dominance give
$E\succeq\delta I$, hence
\[
 \frac{A}{d}
 =2^{-2s}LL^\top+\frac{E}{d\,2^{2s}}
 \succeq\frac{\delta}{d\,2^{2s}}I\succ0.
\]
The proposed $L$ may be obtained numerically, but every equality
and inequality in this acceptance test uses exact integers.
The rational factorization coefficients of an SoS itself are
not required: a positive semidefinite rational Gram admits the
real factorization used in Section~\ref{sec:exact-certificates}.

Run \texttt{python3 verification/positive.py} after installing
the dependencies in \path{verification/requirements.txt}.
The checker validates the complete chain, including the base case,
coefficient identities, central squares, kernel relations, and
positivity. It also rejects deliberately corrupted identities
and Gram corrections. The expected final line is
\begin{quote}\small\ttfamily
PASS: F\_N is SOS for 2 <= N <= 50; 48 exact steps; negative controls rejected
\end{quote}
The compressed coefficients are included with the source archive;
the mathematical acceptance conditions are
\eqref{eq:appendix-normalized-step}--\eqref{eq:exact-residual-certificate}.

\subsection{Lean verification of the large-order negative result}
\label{app:lean-negative}

\rev{The public repository also contains a Lean~4 formalization of the
large-order negative theorem for the original Toeplitz BW polynomial.
In particular, the declaration
\texttt{ToeplitzSOS.Negative.sharpNegativeResolution} proves that the
original real Toeplitz BW form is not SoS for every
$N\ge\SufficientOrder$, where the SoS predicate allows an arbitrary
finite number of homogeneous quadratic squares with arbitrary real
coefficients.  The source is
\path{lean/ToeplitzSOS/Negative/Resolution.lean}.  The verified
formalization is frozen at repository tag \texttt{v1.0}; reproduction
commands and the exact verified source tree are recorded in
\path{VERIFICATION.md}.}

\subsection{The positive Lean statement}\label{app:positive-lean}

In the supplied formalization, \texttt{toeplitzBW n} is the
polynomial defined by the original real Toeplitz matrices, and
\texttt{IsSumSqHomQuad} asserts the existence of a finite sum of
real homogeneous quadratic squares. In the namespace
\texttt{ToeplitzSOS}, the range theorem is
\begin{lstlisting}[basicstyle=\ttfamily\small,breaklines=true]
theorem toeplitzBW_isSumSq_of_le_20
    (n : Nat) (h2 : 2 <= n) (h20 : n <= 20) :
    IsSumSqHomQuad (toeplitzBW n)
\end{lstlisting}
Its proof combines the explicit order-two identity with the
fixed-order theorems for $n=3,\ldots,20$.
The source file in the public repository is
\path{lean/ToeplitzSOS/Certificates/Range.lean}.
The Lean and Mathlib versions are pinned by the repository release.
From the \path{lean/} directory, the positive certificates can be
checked with \texttt{./verify --certificates}.  The rational chain at
orders $21$--$50$ is not part of this theorem.
